\chaptertitle{第三章\quad 数据的集中趋势}

我们经常听到这样的说法："全班数学平均分85分"、"北京居民人均可支配收入X元"、"大多数人的选择是……"。这些说法都用\textbf{一个数}来概括\textbf{一组数据}的特征——它告诉我们这组数据的"中心"大概在哪里。这一章我们学习三种最常用的"代表数"：平均数、中位数和众数。

\sectiontitle{3.1\quad 平均数——最常用的代表数}

平均数是最常见的统计量。它的朴素理解就是"把总数平均分"。

\begin{example}
小明五次数学测验成绩：85、90、92、88、95。求平均分。
\end{example}
\begin{solution}
先把五次成绩加起来：$85+90+92+88+95=450$。再除以次数$5$：$450\div5=90$分。
\end{solution}

但生活中很多情况不是简单的"平均分"——每件事的重要性可能不同。

\begin{example}
小明平时测验平均$90$分（权重$1$），期中考试$88$分（权重$2$），期末考试$94$分（权重$3$）。三种成绩的权重不同，求综合平均分。
\end{example}
\begin{solution}
加权总$=90\times1+88\times2+94\times3=90+176+282=548$，\\
总权重$=1+2+3=6$，加权平均分$=548\div6\approx91.3$分。\\
期末考试占比最高（权重$3$），所以对平均分的影响最大。
\end{solution}

\knowbox{
\textbf{算术平均数}：$\bar x=\dfrac{x_1+x_2+\cdots+x_n}{n}$。\\
\textbf{加权平均数}：$\bar x=\dfrac{w_1x_1+w_2x_2+\cdots+w_nx_n}{w_1+w_2+\cdots+w_n}$，\\
其中$w_i$是第$i$个数据的权重（重要程度）。}

来看一个分组数据求平均的例子——这也是最常见的"坑"之一。

\begin{example}
某班数学成绩分三档：90分以上8人（平均94分），80~89分15人（平均84分），70~79分10人（平均74分）。求全班平均分。
\end{example}
\begin{solution}
千万不能直接用$(94+84+74)\div3$！因为各档人数不同。\\
正确做法：总分数$=8\times94+15\times84+10\times74=752+1260+740=2752$，\\
总人数$=8+15+10=33$，平均分$=2752\div33\approx83.4$分。\\
这里"人数"就是各档成绩的权重。
\end{solution}

\begin{example}
某商店三种糖果单价：奶糖20元/kg、水果糖25元/kg、巧克力糖30元/kg。\\
(1) 三种各买1kg，平均单价是多少？\\
(2) 奶糖买2kg、水果糖买3kg、巧克力糖买5kg，平均单价呢？
\end{example}
\begin{solution}
(1) 重量相同：$(20+25+30)\div3=25$元/kg。\\
(2) 重量不同：总价$=20\times2+25\times3+30\times5=40+75+150=265$元，\\
总重$=2+3+5=10$kg，平均单价$=265\div10=26.5$元/kg。\\
第二种情况就是加权平均——重量是权重。
\end{solution}

\begin{exercise}
\item 求数据$12,15,18,21,14$的平均数。
\item 一组数据的平均数是$10$，原来有$4$个数，加入一个新数$x$后平均数变成$12$，求$x$。
\item 某课程平时成绩占$60\%$、期末成绩占$40\%$。小明平时得85分，期末得90分，求总评。
\item 为什么分组求平均时不能直接把各组的平均数再取平均？
\end{exercise}

\sectiontitle{3.2\quad 中位数——排在中间的那个数}

平均数有一个明显的弱点：它非常容易受极端值的影响。

\begin{example}
某公司$5$名员工的月薪（元）：3000, 3200, 3500, 3800, 50000（总经理）。求平均月薪。
\end{example}
\begin{solution}
平均月薪$=(3000+3200+3500+3800+50000)\div5=12700$元。\\
但前面$4$名员工的月薪都不到$4000$元——"平均月薪12700元"完全不能反映大多数人的收入水平。问题就出在总经理的$50000$元把平均数拉高了。这个例子说明：当数据中有极端值时，平均数的代表性会大打折扣。
\end{solution}

于是我们引入\textbf{中位数}：把数据从小到大排列，取最中间的那个数。如果数据个数是奇数，中位数就是中间那个数；如果是偶数，就是中间两个数的平均数。

\begin{example}
上例中$5$名员工的月薪排序为：3000, 3200, 3500, 3800, 50000。求中位数。
\end{example}
\begin{solution}
奇数个数据（$n=5$），最中间的是第$3$个数：$3500$元。\\
中位数$3500$元不受总经理$50000$元的影响，比$12700$元更能代表"普通员工的收入水平"。
\end{solution}

\begin{example}
$6$名学生的身高（cm）：152, 155, 158, 162, 165, 170。求中位数。
\end{example}
\begin{solution}
偶数个数据（$n=6$），中间两个是第$3$个$158$和第$4$个$162$，\\
中位数$=(158+162)\div2=160$cm。
\end{solution}

\begin{example}
数据$3,7,5,9,11$的中位数是多少？数据$4,8,6,10,12,14$呢？
\end{example}
\begin{solution}
第一组排序后$3,5,7,9,11$，$n=5$奇数，中位数$=7$。\\
第二组排序后$4,6,8,10,12,14$，$n=6$偶数，中位数$=(8+10)\div2=9$。
\end{solution}

\knowbox{
\textbf{中位数}：将数据从小到大排列，\\
若$n$为奇数，取第$\dfrac{n+1}{2}$个数；\\
若$n$为偶数，取第$\dfrac{n}{2}$和第$\dfrac{n}{2}+1$个数的平均数。\\
中位数的优点：不受极端值影响，更能代表"典型"水平。}

\begin{exercise}
\item 求数据$1,3,5,7,9,11$的中位数。
\item 一组数据的中位数是$15$，加入$5$和$25$后，新的中位数可能是多少？说明理由。
\item 在$5$次测验中，小明的成绩分别为$72,85,90,68,95$。求中位数。如果其中一次错记为$100$分（实际是$70$分），对平均数影响大还是对中位数影响大？
\end{exercise}

\sectiontitle{3.3\quad 众数——出现最多的那个数}

有时候我们关心的是"最常见的值"——比如鞋店进货时，更关心哪个尺码卖得最多，而不是平均尺码。

\begin{example}
某鞋店一周内售出的鞋码记录：\\
38, 39, 38, 40, 39, 38, 41, 39, 38, 40, 39, 38, 42, 38, 39。\\
哪种鞋码卖得最多？
\end{example}
\begin{solution}
数一下：38码出现6次，39码出现5次，40码出现2次，41码1次，42码1次。\\
38码出现次数最多，所以\textbf{众数}是38码。鞋店应该多进38码的鞋。
\end{solution}

\knowbox{
\textbf{众数}：一组数据中出现次数最多的数。\\
注：众数可能不止一个（多个数并列出现最多次），也可能没有（各数次数相同）。\\
众数适用于：了解"最常见"的情况（如服装尺码、最喜欢的颜色）。}

\begin{example}
判断以下各组有没有众数，有几个？\\
(1) $5,5,6,7,7,8$ \qquad (2) $1,2,3,4,5$ \qquad (3) $3,3,4,4,5,5$
\end{example}
\begin{solution}
(1) $5$和$7$各出现2次（并列最多），所以有两个众数：$5$和$7$。\\
(2) 每个数只出现1次，没有众数。\\
(3) $3,4,5$各出现2次，有三个众数。
\end{solution}

\begin{exercise}
\item 求数据$2,4,4,5,6,6,6,8$的众数。
\item "众数一定只有一个"这句话对吗？如果不对，举一个反例。
\item 一组数据的众数是$10$，中位数是$12$，平均数是$11$。写出一个满足条件的$5$个数的数据集。
\end{exercise}

\sectiontitle{3.4\quad 三种代表数的比较与选择}

三种代表数各有所长，也各有所短。我们用一张表来对比：

\begin{center}
\begin{tabular}{c|c|c|c}
 & 平均数 & 中位数 & 众数\\\hline
怎么算 & 总和$\div$个数 & 排序取中间 & 看谁出现最多\\
受极端值影响 & 很大 & 不敏感 & 不敏感\\
适用范围 & 数据分布均匀 & 有极端值时 & 需要"常见值"时\\
例子 & 平均身高 & 房价、收入 & 鞋码进货\\
局限 & 怕"拖后腿"的数据 & 只用了排序信息 & 可能没有或不唯一
\end{tabular}
\end{center}

\begin{example}
某公司招聘，广告说"本公司员工平均月薪$8000$元"。你入职后发现普通员工月薪只有$4000$元，高管月薪$3$万元。以下哪个说法正确？\\
A. 广告说谎了；B. 平均数没骗人，但中位数更能代表普通员工；C. 应该看众数。
\end{example}
\begin{solution}
选B。平均数确实可能是$8000$——如果高管足够多，但$4000$才是大多数人的实际收入。此时中位数（或者最低的月薪）比平均数更能反映真实情况。这就是平均数被极端值"拉高"的典型案例。
\end{solution}

\begin{example}
班主任想了解班里大多数同学的数学成绩在哪个分数段。她算了平均分83分，又看了中位数85分，最后发现85分出现了8次（众数=85）。三个数相差不大，说明这次考试的成绩分布比较均匀。但如果平均分70、中位数80、众数90——那就说明有少数极低分拉低了平均分，而多数人考得不错。
\end{example}

\begin{exercise}
\item 以下情况用平均数、中位数还是众数更合适？\\
(a) 鞋店决定下个月多进哪个尺码的鞋；\\
(b) 比较两个班的数学成绩；\\
(c) 报告一个城市的房价水平让人了解"普通人买得起什么价位的房子"。
\item 某公司$20$人，$19$人的月薪在$4000$~$6000$元，CEO月薪$10$万元。平均月薪被拉高到约$10000$元。如果你是该公司的普通员工，你觉得"平均月薪"能代表你的收入水平吗？
\end{exercise}

\sectiontitle{本章小结}

这一章我们学习了三个描述"数据集中趋势"的统计量。平均数是最常用的，它利用了全部数据，但害怕极端值——一个特别大或特别小的数就能把它拉偏。中位数不怕极端值，因为它只关心排序后中间位置的值，但这也意味着它忽略了其他数据的信息。众数告诉我们"最常见的是什么"，对于分类数据（如最喜欢的颜色、最畅销的鞋码）特别有用。在实际分析中，常常把三者结合起来看——它们从不同角度告诉我们不同的信息。

\sectiontitle{课后练习}

\subsection*{A组\quad 基础巩固}

\begin{exercise}
\item 求下列各组数据的平均数、中位数和众数：\\
(a) $4,7,9,10,15$ \qquad (b) $2,5,5,6,8,9$ \qquad (c) $10,10,10,10,10$
\item 某小组$8$人数学成绩：$72,85,90,68,95,88,90,76$，求平均分和中位数。
\item 数据$3,7,8,10,x$的中位数是$8$，求$x$的最小整数值。
\item 两组数据：A组$5,5,5,5,5$；B组$3,4,5,6,7$。它们的平均数、中位数、众数各是多少？
\item 如果一组数据的平均数、中位数和众数都相等，这组数据可能有什么特征？
\item 某门课平时成绩占$40\%$、期中占$20\%$、期末占$40\%$。小明平时$88$分、期中$76$分、期末$92$分，求总评成绩。
\end{exercise}

\subsection*{B组\quad 挑战提升}

\begin{exercise}
\item 某班$30$人，其中$29$人的平均身高$155$cm。班主任身高$175$cm，加入后全班平均身高变多少？这个变化说明了什么？
\item 数据$a,b,c,d$的平均数是$10$，中位数是$12$。$a<b<c<d$且都是正整数，求所有可能的取值。
\item 某公司$20$人：普通员工$15$人（年薪$8$万），中层$4$人（$20$万），高管$1$人（$100$万）。\\
(a) 求平均数和中位数。\\
(b) 公司想用哪个统计量来宣传"待遇优厚"？你作为求职者会怎么看？
\item 有$5$个数，平均数是$12$，中位数是$13$，众数是$14$。写出满足条件的一组数。
\item 某班$10$次测验的平均分是$82$分，其中$2$次特别难（平均$50$分）。去掉这$2$次，剩下$8$次的平均分是$90$分。为什么包含低分时的平均数$82$分不能反映班级的正常水平？如果可以，你会怎么改进"平均分"的计算方式？
\end{exercise}

\newpage
